\chapter{Repères de lecture des variables}
\label{ann:variables}

La liste décrit les principaux champs du rapport. Elle ne remplace pas le dictionnaire
 de l'ONISR ni le schéma des 119 colonnes préparées. Les codes catégoriels restent des
étiquettes. Seuls les champs retenus au chapitre~\ref{chap:protocole-ml} alimentent le modèle.

{\small\renewcommand{\arraystretch}{1.18}
\begin{longtable}{@{}p{0.34\textwidth}p{0.61\textwidth}@{}}
\caption{Variables et indicateurs utiles à la lecture du rapport.}\label{tab:dictionnaire}\\
\toprule Champ & Signification et précaution\\\midrule
\endfirsthead
\toprule Champ & Signification et précaution (suite)\\\midrule
\endhead
\midrule\multicolumn{2}{r}{Suite à la page suivante}\\
\endfoot
\bottomrule\endlastfoot
\code{Num_Acc} & Identifiant d'accident, partagé par ses usagers.\\
\code{id_vehicule} & Identifiant du véhicule ; utilisé avec l'accident pour la jointure.\\
\code{id_usager} & Identifiant officiel lorsqu'il existe, absent en 2019-2020.\\
\code{id_observation} & Clé technique année + rang source, sans suivi longitudinal.\\
\code{grav} & Gravité préparée : indemne, tué, hospitalisé ou léger.\\
\path{cible_deces_hospitalisation} & Cible 1 pour décès ou hospitalisation, 0 pour indemne ou léger ; inconnue conservée comme absence.\\
\code{cible_connue} & Indique si la gravité permet de définir la cible binaire.\\
\code{catu} & Catégorie : conducteur, passager ou piéton.\\
\code{age_approx} & Année de l'accident moins naissance, contrôlée entre 0 et 110 ans.\\
\code{tranche_age} & Variable créée pour le modèle ; huit tranches et une modalité d'absence.\\
\code{sexe} & Sexe enregistré dans BAAC, avec modalité d'absence pour le modèle.\\
\code{catv} & Catégorie du véhicule associé ; véhicule heurtant pour un piéton.\\
\code{lum} & Conditions de lumière ; cinq modalités connues.\\
\code{agg} & Hors agglomération ou en agglomération.\\
\code{int} & Type d'intersection.\\
\code{atm} & Conditions atmosphériques enregistrées.\\
\code{mois_accident} & Mois de la date de l'accident, traité comme catégorie.\\
\code{jour_semaine} & Jour de semaine dérivé de la date, traité comme catégorie.\\
\code{heure_accident} & Heure extraite de l'horaire, traitée comme catégorie.\\
\code{lieu_vma} & Limitation autorisée résumée au niveau de l'accident ; hors modèle ML.\\
\code{lieu_vma_desaccord} & Plusieurs valeurs connues de limitation pour l'accident.\\
\code{lieu_vma_partiel} & Valeur connue coexistante avec une valeur inexploitable.\\
\code{lieu_vma_a_verifier} & Au moins une limitation source préparée supérieure à 130.\\
\code{age_manquant} & Âge préparé absent, y compris après contrôle.\\
\code{age_suspect} & Âge brut hors de l'intervalle retenu.\\
\code{role_temporel} & Apprentissage, validation ou test selon l'année.\\
\end{longtable}}

\FloatBarrier
